\subsection{Compte-Titres Ordinaire (CTO)}
\label{sec:cto}

Le CTO, contrairement à l'assurance-vie, ne bénéficie d'aucun régime fiscal spécifique en matière successorale, le CTO intègre donc l'ensemble du patrimoine du défunt. Toutefois, il présente une caractéristique importante : l'effacement des plus-values latentes au décès, ce qui confère un avantage aux héritiers \cite{lechiquier_2019}.

\subsubsection{Fonctionnement et frais}
\label{sec:cto_evolution_capital}

Pour comparer rigoureusement les deux enveloppes, nous supposons que l'épargnant investit dans les mêmes supports financiers (par exemple, un ETF suivant un indice mondial). Dans ce cas, les frais internes aux supports sont identiques.
La différence majeure est que le CTO n'a \textbf{pas de frais de gestion annuels} propres à l'enveloppe. À rendement brut égal, le capital sur un CTO croît donc plus vite que sur une assurance-vie, car il n'est pas amputé chaque année par des frais d'enveloppe.

\subsubsection{Fiscalité en cas de transmission}
\label{sec:cto_base_taxable}

Au décès, le CTO présente deux caractéristiques fiscales majeures :
\begin{enumerate}
    \item \textbf{La purge des plus-values latentes :} C'est l'avantage principal du CTO. Les héritiers reçoivent les titres valorisés au jour du décès. Toute la plus-value accumulée du vivant du souscripteur est \textbf{définitivement exonérée d'impôt sur le revenu et de prélèvements sociaux}. Cet avantage est considérable, surtout après une longue période de détention et de forte croissance.
    \item \textbf{L'intégration à la succession classique :} La valeur totale du CTO au jour du décès est ajoutée aux autres biens du défunt (immobilier, liquidités, etc.) pour former l'actif successoral. Cet actif est ensuite soumis au barème des droits de succession classique, après application des abattements légaux (par exemple, 100 000 € par enfant en ligne directe).
\end{enumerate}

Le barème progressif en ligne directe s'applique par tranches (5\%, 10\%, 15\%, 20\%, etc.) \cite{bofip_droits_mutation_general}. L'héritage net correspond à la valeur du CTO, diminuée de la part des droits de succession qui lui est imputable.
